% Packages
\usepackage{amsmath, mathtools, bm, commath, amssymb}
\usepackage{tensor}
\usepackage{physics}
\usepackage{siunitx}
\usepackage{nicematrix}

% General settings
\allowdisplaybreaks     % Allows for equations to break between pages

% Quick constants for lazyness
\newcommand{\pihalf}{\frac{\pi}{2}}
\newcommand{\pithird}{\frac{\pi}{3}}
\newcommand{\pifourth}{\frac{\pi}{4}}

% Cancel-lines related
\usepackage[thicklines]{cancel}
\renewcommand\CancelColor{\color{xred}}
\newcommand{\cancelcol}[2][xred]{ % This is such a silly solution...
	\renewcommand\CancelColor{\color{#1}}
	\cancel{#2}
	\renewcommand\CancelColor{\color{xred}}
}

% Easier space-notation
\newcommand{\Rs}[1][]{\mathbb{R}^{#1}}
\newcommand{\Cs}[1][]{\mathbb{C}^{#1}}

% Conjugation
\newcommand{\conj}[1]{\overline{#1}}

% Point labels
\newcommand{\pnt}[1]{\mathbf{#1}}

% Notation for vectors (currently: bold letters)
% \renewcommand{\vec}[1]{\bm{#1}}
\newcommand{\uvec}[1]{\hat{#1}}
\newcommand{\vnorm}[1]{\left\| \vec{#1} \right\|}
\newcommand{\cvec}[1]{\bm{\overline{#1}}}
\newcommand{\mat}[1]{\mathbf{#1}}

% Identity matrix
\newcommand{\Id}[1]{\mathbb{I}_{#1}}

% Vector operations
\newcommand{\inner}[2]{\langle #1,#2 \rangle}

% Dual space related
\newcommand{\dualspace}[1]{#1^{*}}
\newcommand{\dualvec}[1]{\vec{#1}^{*}}
\newcommand{\dualbasis}[1]{#1^{*}}
\newcommand{\dualeb}[1]{\vec{\epsilon}^{#1}}
\newcommand{\dualebc}[1]{\vec{\epsilon}_{#1}}
\newcommand{\dualRs}[1][]{\mathbb{R}^{#1*}}
\newcommand{\dualCs}[1][]{\mathbb{C}^{#1*}}

% Row- and column-vectors: arguments separated by ";".
% Example: $\vec{a} = \colvec{1;2;3;4}$.
\makeatletter
\newcommand\rcvector[2][\\]{\ensuremath{%
		\global\def\rc@delim{#1}%
		\negthinspace\begin{bmatrix}
			\rc@vector #2;\relax\noexpand\@eolst%
		\end{bmatrix}}}
\def\rc@vector #1;#2\@eolst{%
	\ifx\relax#2\relax
		#1
	\else
		#1\rc@delim
		\rc@vector #2\@eolst%
	\fi}
\makeatother
\newcommand{\colvec}{\rcvector}
\newcommand{\rowvec}[1]{\rcvector[\ \;]{#1}}
\newcommand{\GenericRowVec}[2][n]{\rowvec{#2_{1};#2_{2};\dots;#2_{#1}}}
\newcommand{\GenericColVec}[2][n]{\colvec{#2^{1};#2^{2};\vdots;#2^{#1}}}

% Colored vectors
\newcommand{\colorVec}[2]{\color{#1}{\vec{#2}}\color{black}}
\newcommand{\vred}{\colorVec{xdarkred}{v}}
\newcommand{\vblue}{\colorVec{xdarkblue}{v}}
\newcommand{\vgreen}{\colorVec{xdarkgreen}{v}}

% Basis vectors and transformations
\newcommand{\eb}[1]{\hat{e}_{#1}}
\setlength{\fboxsep}{1pt} % This makes the colored e1, e2 and e3 more compact
\newcommand{\ebx}{\colorbox{xred!25}{$\eb{1}$}}
\newcommand{\eby}{\colorbox{xblue!25}{$\eb{2}$}}
\newcommand{\ebz}{\colorbox{xgreen!25}{$\eb{3}$}}

\newcommand{\ebc}[1]{\tilde{\vec{e}}_{#1}}
\newcommand{\ebr}[1]{\color{xdarkblue}\eb{#1}\color{black}}
\newcommand{\ebcr}[1]{\color{xdarkred}\ebc{#1}\color{black}}
\newcommand{\oldB}{\color{xdarkblue}B\color{black}}
\newcommand{\newB}{\color{xdarkred}\tilde{B}\color{black}}
\newcommand{\Forw}{\bm{F}}
\newcommand{\Backw}{\bm{F^{-1}}}

% Transformation matrices components
\newcommand{\Forwcomps}{F\indices{^i_j}}
\newcommand{\Backwcomps}{B\indices{^i_j}}

% Dual versions
\newcommand{\dualOldB}{\color{xdarkblue}\dualbasis{B}\color{black}}
\newcommand{\dualNewB}{\color{xdarkred}\tilde{B}^{*}\color{black}}

% Dark equal?
\newcommand{\beq}{\color{black}{=}}

% Better imaginary unit and natural base notation (to separate from variables)
\newcommand{\iu}{\mathrm{i}\mkern1mu}
\newcommand{\eu}{\mathrm{e}}
\newcommand{\Eu}[1]{\mathrm{e}^{#1}}
\newcommand{\EX}[1]{\exp\left(#1\right)}

% General nice matrix
\newcommand{\GNMatrix}[3]{
	\begin{bNiceMatrix}
		#1_{11}  & #1_{12}  & \dots  & #1_{1#3}  \\
		#1_{21}  & #1_{22}  & \dots  & #1_{2#3}  \\
		\vdots   & \vdots   & \Ddots & \vdots    \\
		#1_{#21} & #1_{#22} & \dots  & #1_{#2#3}
	\end{bNiceMatrix}
}

% Color-coded linear combinations
\newcommand{\LCa}[1]{\textcolor{xred}{#1}}
\newcommand{\LCb}[1]{\textcolor{xblue}{#1}}
\newcommand{\LCc}[1]{\textcolor{xdarkgreen}{#1}}
\newcommand{\LCd}[1]{\textcolor{xpurple}{#1}}
\newcommand{\LCe}[1]{\textcolor{xorange}{#1}}

% Colored column vectors? Should be implemented better...
% (but I need to write instead of over-engineering the code)
\newcommand{\clrcolvecxy}[2]{
	\colvec{\LCa{#1};\LCb{#2}}
}
\newcommand{\clrgenericvec}[2]{
	\colvec{\LCa{#1_{1}};\LCb{#1_{2}};\vdots;\LCe{#1_{#2}}}
}

\newcommand{\spacefortrans}[1]{
	\begin{axis}[
			width=7cm,
			height=7cm,
			xyplane,
			xmin=-5, xmax=5,
			ymin=-5, ymax=5,
			enlargelimits={abs=0.0},
			grid=major,
			major grid style={black!50},
			xtick={-5,...,5},
			ytick={-5,...,5},
			anchor=origin,
			xlabel={},
			ylabel={},
			domain={-5:5},
			#1,
		]
		% Rectangles
		\draw[very thick, xorange, fill=xorange!25] (0,0) rectangle (1,1);
		\draw[very thick, xblue, fill=xblue, fill opacity=0.25] (2,2) rectangle (4,3);

		% Basis vectors
		\draw[vector, xred] (0,0) -- (1,0);
		\draw[vector, xblue] (0,0) -- (0,1);

		% Parallel and orthogonal lines
		\draw[very thick, densely dotted, xpurple] (-5,-3) -- (5,-1);
		\draw[very thick, densely dotted, xpurple] (-5,-4) -- (5,-2);

		% Smooth shape
		\draw[very thick, xgreen, fill=xgreen, fill opacity=0.25] (-2.5,1) to[closed, curve through = {(-3.5,-1) (-4.5,0) (-3.5,1) (-3.5,3) (-2.5,2)}] (-0.5,2); % This shape is shamelessly stolen from https://tikz.org/examples/chapter-12-drawing-smooth-curves/

		% A small circle at the origin
		\draw[xpurple, very thick, fill=white] (0,0) circle (2pt);
	\end{axis}
}

\newcommand{\drawtrans}[7]{
	% This function needs to be re-defined in terms of
	% NewDocumentCommand or whatever it's called
	\begin{tikzpicture}
		\spacefortrans{name=beforetrans, xlabel={$x$}, ylabel={$y$}}
		\pgftransformcm{#2}{#3}{#4}{#5}{\pgfpoint{#6}{#7-10cm}}
		\spacefortrans{name=aftertrans}
		\pgftransformreset
		\draw[|->, very thick] ($(beforetrans.south)+(0,-8mm)$) -- ++(0,-2cm) node[midway, right] {\Huge$#1$};
		\node[anchor=west] at (aftertrans.east) {$x'$};
		\node[anchor=south] at (aftertrans.north) {$y'$};
	\end{tikzpicture}
}

\newcommand{\fromtranstomat}[6]{
	\begin{tikzpicture}[ampersand replacement=\&]
		\tikzset{mapstoarrow/.style={|->, thick, xblue}}

		\draw[thick, dotted, black!25] (-2,-2.25) -- ++(4,0) node[midway, fill=white, text=xblue] (T) {$#1$};

		\node (e1) at (-1.,0) {$\eb{1}$};
		\node (e2) at (+1.,0) {$\eb{2}$};
		\node[below of=e1] (e1col) {$\colvec{1;0}$};
		\node[below of=e2] (e2col) {$\colvec{0;1}$};
		\node[below of=e1col, yshift=-1.2cm] (Te1) {$#1\left(\eb{1}\right)$};
		\node[below of=e2col, yshift=-1.2cm] (Te2) {$#1\left(\eb{2}\right)$};
		\node[below of=Te1] (e1map) {$\colvec{#3;#4}$};
		\node[below of=Te2] (e2map) {$\colvec{#5;#6}$};
		\draw[mapstoarrow] (e1col.south) -- (Te1.north);
		\draw[mapstoarrow] (e2col.south) -- (Te2.north);

		\matrix [
		matrix of math nodes,
		nodes={rectangle, minimum size=15pt},
		inner sep=0pt,
		column sep=8pt,
		left delimiter={[}, right delimiter={]},
			] (A) at (0,-6.5)
		{
				#3 \& #5 \\
				#4 \& #6 \\
			};
		\begin{scope}[on background layer]
			\filldraw[xpurple, rounded corners, opacity=0.25] (A-1-1.north west) rectangle (A-2-1.south east);
			\filldraw[xorange, rounded corners, opacity=0.25] (A-1-2.north west) rectangle (A-2-2.south east);
		\end{scope}

		\draw[-stealth, thick, xpurple] (e1map.south) to[out=270, in=90] (A-1-1.north);
		\draw[-stealth, thick, xorange] (e2map.south) to[out=270, in=90] (A-1-2.north);
		\node[below of=A, text=xblue] {$#2$};
	\end{tikzpicture}
}
